\section{Validation Scope and Context}

\subsection{Intent of the Campaign}
This report is meant for validation activities where the key question is whether
the design satisfies documented requirements, regulatory expectations, or
internal release criteria. It should read as an argument supported by evidence,
not just a loose collection of bench notes.

\specbox{
Validation reports are strongest when they connect requirements, evidence,
exceptions, and sign-off in one place. Reviewers should not need a second file
to understand whether the product is ready.
}

\subsection{Boundary Conditions}
\begin{table}[H]
    \centering
    \small
    \begin{tabular}{l p{9.5cm}}
        \toprule
        \textbf{Area} & \textbf{Validation Boundary} \\
        \midrule
        Functional Scope & Features, modes, or mission states covered by this campaign. \\
        Exclusions & Known out-of-scope conditions, deferred scenarios, or pending design changes. \\
        Inputs Considered & Supply range, sensor domain, communication states, or external commands. \\
        Outputs Evaluated & Performance, compliance, protection behavior, user-visible reactions, or logs. \\
        Evidence Sources & Bench data, simulation correlation, analysis, review records, or supplier documentation. \\
        \bottomrule
    \end{tabular}
    \caption{Suggested validation boundary statement}
    \label{tab:validation-boundary}
\end{table}

\subsection{Approval Expectations}
\makeapprovaltable{
    Validation Owner & Reviewer Name & Pending & --- \\
    Systems Engineering & Reviewer Name & Pending & --- \\
    Product / Project Lead & Reviewer Name & Pending & --- \\
}
